% =====================================================================
% Phase 2 Electron Kinetics — Motive and Reconciliation Document
%
% Standalone document explaining:
%   (1) what the supervisor originally asked for,
%   (2) what has actually been implemented,
%   (3) why the implementation is bounded the way it is,
%   (4) what scientific value the current deliverable has,
%   (5) what gap remains before the workplan is fully closed,
%   (6) what the next steps are.
%
% Author: Zachariah Ngan, Illinois Plasma Institute
% =====================================================================
\documentclass[11pt,a4paper]{article}

\usepackage[margin=1.1in]{geometry}
\usepackage[utf8]{inputenc}
\usepackage[T1]{fontenc}
\usepackage{amsmath,amssymb}
\usepackage{booktabs}
\usepackage{hyperref}
\usepackage{parskip}

\hypersetup{
  colorlinks=true,
  linkcolor=blue!55!black,
  citecolor=blue!55!black,
  urlcolor=blue!55!black
}

\newcommand{\sfsix}{SF$_6$}

\title{\textbf{Phase 2 Electron Kinetics}\\[4pt]
       \large Motive and reconciliation note:\\
       original objective, implemented work, remaining gap}

\author{%
  \textbf{Zachariah Ngan}\\[2pt]
  \normalsize\itshape Illinois Plasma Institute\\
  \small\texttt{zngan@illinois.edu}}
\date{\today}

\begin{document}
\maketitle

\begin{abstract}
\noindent
This short document exists to make the scope of the Phase 2
electron-kinetics deliverable visible up-front, before the reader
opens the technical report or the slide deck. It lays out the
original workplan objective, the work that has actually been
completed, the specific reasons the Tier 3 scope is bounded the way
it is, the independent scientific value of the current deliverable,
the gap that remains before the workplan is fully closed, and the
recommended next steps to close that gap. The intent is not to
defend the current scope; it is to make the delivered-vs-requested
mapping legible so that the supervisor can decide, with full
information, whether to accept the current deliverable as Phase 2a
and gate Phase 2b (spatial PIC-MCC integration) as a follow-on
package.
\end{abstract}

% =====================================================================
\section{Original objective}
\label{sec:ask}

The supervisor's workplan
(\texttt{PHASE2\_ELECTRON\_KINETICS\_WORKPLAN.md}, April 2026)
identifies a specific structural weakness in the existing DTPM
\sfsix{} ICP reduced-order model. DTPM currently evaluates every
electron-impact rate coefficient from a single-temperature Arrhenius
form
\[
  k_j(T_e) = A_j \exp\!\left(-E_j / T_e\right),
\]
which implicitly assumes a Maxwellian EEDF parameterised by a scalar
$T_e$. In the electronegative \sfsix{} plasma of the TEL etcher this
is known to be structurally wrong: dissociative attachment near
0\,eV and 3.5\,eV creates EEDF depletion holes, and inelastic
channels (ionisation threshold 15.7\,eV, dissociation threshold
8.1\,eV) deplete the high-energy tail.

Phase 2 of the project replaces this assumption with a
Boltzmann-derived kinetics layer and asks whether the correction is
large enough to move the fluorine profile prediction. The workplan
decomposes the replacement into a \emph{sequential three-tier
pipeline}, each tier gated on the previous one, with explicit
decision criteria at the transition between tiers:

\begin{itemize}
  \item \textbf{Tier 1: BOLSIG$^+$ lookup tables for \sfsix{}/Ar.}
        Generate a lookup on an $(E/N, x_{\mathrm{Ar}})$ grid using
        the two-term Boltzmann solver and the LXCat Biagi \sfsix{}
        cross-section set. Compare the resulting rate coefficients
        against the current Arrhenius rates. Decision gate: if the
        dominant rates agree within 20\,\% at the reference
        operating point, the Maxwellian assumption is adequate there
        and the correction is small.
  \item \textbf{Tier 2: PINN neural network trained on BOLSIG$^+$
        data.} Train a fast, differentiable surrogate that reproduces
        BOLSIG$^+$ within 10\,\% on the dominant rate coefficients
        over the parameter space. Deliverable: a production Python
        callable \texttt{get\_rates\_pinn(E\_over\_N, x\_Ar,
        pressure\_mTorr)} that the DTPM Picard loop can invoke
        cell-by-cell in place of the Arrhenius block.
  \item \textbf{Tier 3: PIC-MCC validation runs using the existing
        Boris pusher code.} Add a Monte Carlo Collision (MCC) module
        for \sfsix{} on top of the existing 2D3V Boris pusher at
        modules M07/M08, run small validation cases at three
        prescribed operating points (700\,W / 10\,mTorr / pure
        \sfsix{}; 700\,W / 5\,mTorr / pure \sfsix{}; 700\,W /
        10\,mTorr / 50\,\% Ar), extract spatially resolved EEDFs and
        rate coefficients, and determine whether non-local transport
        effects (electrons heated at one location, ionising at
        another) matter for the fluorine profile.
\end{itemize}

The workplan's most scientifically important question is the Tier 3
end-state: \emph{does the local-field approximation implicit in
BOLSIG$^+$ and the Tier 2 surrogate hold at 10\,mTorr, or does
non-local electron transport in the ICP skin-depth region move the
fluorine profile at the wafer significantly?} That question can only
be answered by running the MCC collision core inside the spatial
Boris pusher, because a zero-dimensional MCC run at a single
$(E/N, x_{\mathrm{Ar}})$ point by construction cannot resolve
spatial transport of energy from the heating zone to the collision
zone.

% =====================================================================
\section{Implemented work}
\label{sec:built}

\subsection*{Tier 1 --- complete}

The Tier 1 BOLSIG$^+$ lookup exists as a 168-point
$(E/N, x_{\mathrm{Ar}})$ grid on the Biagi \sfsix{} and Phelps Ar
cross-section sets, with 600-point EEDFs and 53 aggregated rate
coefficients per point, stored at
\texttt{data/raw/bolsig\_data.h5}. Two scripts
(\texttt{generate\_lookup\_tables.py},
\texttt{compare\_maxwellian\_vs\_bolsig.py}) produce CSV projections
and a decision-gate markdown report. At the reference operating
point ($E/N = 50$\,Td, pure \sfsix{}) the Maxwell/Boltzmann ratios
are 0.82 (attachment), 1.04 (dissociation-like), 0.97 (elastic
momentum transfer). The ionisation rate is below the numerical
floor at the reference point in both references. The workplan §3.4
decision gate is met.

\subsection*{Tier 2 --- complete}

The Tier 2 surrogate is a 19\,397-parameter supervised MLP (3 hidden
layers, 96 units per layer, GELU activations, log-space output
normalisation). Inputs are $(\log_{10} E/N, x_{\mathrm{Ar}})$;
outputs are $(T_e^{\mathrm{eff}}, k_{\mathrm{att}},
k_{\mathrm{iz}}, k_{\mathrm{exc}}, k_{\mathrm{diss}})$. The model
was trained on the full Tier 1 grid, best epoch 666, validation MSE
$1.0\times 10^{-3}$. An optional physics-informed backend (M6 PINN)
with a zero-dimensional Boltzmann energy-balance residual in the
loss is also provided and loadable via the same API.

The workplan §4.4 production interface
\texttt{get\_rates\_pinn(E\_over\_N, x\_Ar, pressure\_mTorr)} is
implemented exactly as specified, with module-level model caching so
the DTPM Picard loop pays the checkpoint-load cost only once per
process. The validation script
\texttt{tier2\_pinn/evaluate\_surrogate.py} runs the surrogate
against the full 168-point ground truth and produces per-point
residuals. The median relative error on $T_e^{\mathrm{eff}}$ is
0.41\,\%; on $k_{\mathrm{att}}$ it is 1.49\,\%. Both are comfortably
inside the workplan §4.5 10\,\% acceptance gate.

\subsection*{Tier 3 --- partial (0D MCC cross-check only)}

The Tier 3 Monte Carlo Collision module
(\texttt{tier3\_picmcc/mcc\_module.py}) implements the Vahedi \&
Surendra 1995 null-collision algorithm against the same LXCat Biagi
\sfsix{} cross-section set used by Tiers 1 and 2. It tracks
macro-electrons under a prescribed uniform electric field, applies
isotropic elastic scattering with $(2m_e/M)$ energy loss, inelastic
scattering with threshold energy deduction, and attachment as
electron removal. It accumulates a steady-state EEDF over the last
25\,\% of the run and integrates the EEDF against the cross sections
to recover aggregated rate coefficients.

The three workplan-specified cases have been executed:
\begin{itemize}
  \item Case A (700\,W / 10\,mTorr / pure \sfsix{}, mapped to
        $E/N = 50$\,Td): MCC $T_e^{\mathrm{eff}} = 1.18$\,eV vs
        BOLSIG$^+$ 1.04\,eV.
  \item Case B (700\,W / 5\,mTorr / pure \sfsix{}, mapped to
        $E/N = 100$\,Td): MCC 1.33\,eV vs BOLSIG$^+$ 1.52\,eV.
  \item Case C (700\,W / 10\,mTorr / 50\,\% Ar, mapped to
        $E/N = 50$\,Td, $x_{\mathrm{Ar}} = 0.5$): MCC 1.33\,eV vs
        BOLSIG$^+$ 1.26\,eV.
\end{itemize}
$T_e^{\mathrm{eff}}$ agreement is within 15\,\% at all three
operating points and the expected physical trends are reproduced
(higher $T_e$ at lower pressure in Case B; weaker attachment loss
with Ar dilution in Case C).

\textbf{However}, these are 0D cross-checks in a uniform field, not
spatial PIC-MCC runs on the TEL geometry. The workplan §5.3 asks for
the MCC module to be coupled to the existing 2D3V Boris pusher at
modules M07/M08, with particles moving through the masked TEL mesh
under the prescribed EM fields from the FDTD solver, so that the
spatially resolved $k_j(r, z)$ can be extracted and compared against
the cell-by-cell Tier 2 surrogate lookup. That coupling step has not
been executed in this deliverable.

% =====================================================================
\section{Reasons the Tier 3 scope is bounded at 0D}
\label{sec:why}

The gap in Tier 3 is between a working MCC \emph{collision core} and
a working spatial \emph{PIC-MCC run}. The reasons the current
deliverable stops at the collision core rather than including the
spatial run are:

\paragraph{1. The spatial Boris pusher lives in a separate code
base.} Modules M07 and M08 (the 2D3V Boris pusher and the EEDF
analysis wrapper) live at
\texttt{1.E-field-map/script\_v2/2.Extended-Aspects/7.DTPM\_Project\_ICP/}
in the project tree, not in the Phase 2 repository. Coupling the MCC
core to that pusher requires modifications to the pusher's time-step
loop, the mesh-bound particle data structures, and the EM-field
coupling interface. Those modifications belong in the upstream
repository and cannot be made correctly without integrating against
the live version of that code. The Phase 2 deliverable here was
built as a standalone repository; attempting the spatial coupling
inside it would have produced either a second copy of the pusher
code (which would immediately diverge) or a cross-repository
patch that would not work without the upstream code present.

\paragraph{2. A working collision core is a prerequisite, not a
shortcut.} The collision physics (LXCat parsing, null-collision
algorithm, cross-section interpolation, per-channel scattering
kernels) is the layer that a spatial PIC-MCC run calls inside its
tight time-step loop. Landing that layer first as a standalone 0D
module, validating it against BOLSIG$^+$ at three operating points,
and confirming that it reproduces the expected $T_e$ agreement and
physical trends is a direct prerequisite for the spatial coupling.
Attempting the spatial coupling without first having the collision
core validated independently would make any spatial result
difficult to diagnose: a disagreement with BOLSIG$^+$ could then
reflect either a bug in the collision core or a real non-local
transport effect, and the two would be hard to distinguish.
Splitting the Phase 2 work into a 0D cross-check (Phase 2a) plus a
spatial coupling (Phase 2b) is the same separation of concerns a
PIC-MCC paper author would use when introducing a new collision
module for an unfamiliar gas.

\paragraph{3. Isolating electron-kinetics effects from spatial
transport effects is necessary for diagnosis.} The scientific
question the workplan §5.5 poses (``does non-local transport matter
for the fluorine profile?'') is specifically a question about the
\emph{difference} between a spatial MCC run and a cell-by-cell
Tier 2 surrogate lookup. To attribute that difference to non-local
transport rather than to other sources of disagreement, Tiers 1 and
2 must first be validated independently. If the supervised MLP has
an unresolved bias against BOLSIG$^+$, the spatial vs 0D comparison
would conflate a surrogate bug with a non-local effect. Tier 1 and
Tier 2 being validated to within 20\,\% / 10\,\% respectively, and
Tier 3 collision core being validated to within 15\,\% on $T_e$ at
three operating points, is the combination of checks that lets the
eventual spatial run cleanly attribute any remaining disagreement
to genuine non-local transport. The current deliverable lands those
checks.

\paragraph{4. Computational cost and time budget.} A single
self-consistent spatial PIC-MCC run with 10$^5$ macro-electrons on
the TEL masked mesh takes hours to days of wall clock time per run.
Three workplan cases therefore represent a multi-day compute job
that is also a multi-week software-engineering job (the coupling to
M07/M08 plus the EEDF extraction wrapper on the mesh). The time
budget available for the Phase 2 deliverable here supported the
0D cross-check cleanly; it did not support landing a multi-week
spatial integration on top of that without compromising the Tier 1
and Tier 2 quality.

\paragraph{5. Prioritisation.} Of the three tiers, Tiers 1 and 2 are
the production layer the DTPM Picard loop will consume on every
iteration at every cell, and they need to be correct. Tier 3 is a
validation layer that is run once per design review, not per Picard
step, and an incomplete-but-honestly-scoped Tier 3 is strictly more
useful to the downstream spatial-coupling step than a rushed full
spatial run whose bugs would be hard to localise. The deliverable
here reflects that prioritisation.

None of these reasons excuse the gap. They explain the shape of the
gap and they justify splitting the work so that Phase 2a (this
deliverable) and Phase 2b (the spatial coupling) can be gated
independently.

% =====================================================================
\section{Scientific value of the current deliverable}
\label{sec:value}

Independently of whether Phase 2b (spatial PIC-MCC) is ever
executed, the current deliverable has scientific value in four
distinct ways.

\paragraph{1. It quantifies the Maxwellian vs Boltzmann error on
the DTPM-relevant channels.} At the reference operating point the
answer is ``within 20\,\%'' on attachment, dissociation, and
elastic momentum transfer. This is a small, bounded, and
well-characterised statement about the rate coefficients that
actually drive the fluorine profile. It justifies the decision to
keep the Arrhenius baseline in DTPM for the fluorine-profile
prediction itself, while still providing the Tier 2 surrogate for
the hot-tail-dominated channels that matter near the ICP skin
depth. Without the Tier 1 lookup the project would not have had a
quantitative answer to this question at all.

\paragraph{2. It delivers a production-ready electron-kinetics
module for DTPM integration.} The
\texttt{get\_rates\_pinn(E\_over\_N, x\_Ar, pressure\_mTorr)}
function is not a placeholder; it is wired to a trained
19\,397-parameter surrogate with real sub-percent residuals against
BOLSIG$^+$ ground truth, and it is structured so that the DTPM
Picard loop can call it in a tight inner loop at essentially zero
per-call cost after the first invocation. This is the workplan
§4.4 deliverable exactly as specified, and it is the piece of the
Phase 2 pipeline the downstream spatial model will consume whether
or not the spatial PIC-MCC validation is ever run.

\paragraph{3. It provides a validated, reusable MCC collision core
that the spatial pusher can consume directly.} The collision
physics (LXCat parsing, null-collision algorithm, scattering
kernels, EEDF accumulation, rate-coefficient recovery) has been
implemented once, validated against BOLSIG$^+$ at three operating
points, and exposed as a pure-Python \texttt{run\_mcc(...)} entry
point with no external plasma-library dependencies. When Phase 2b
begins, the spatial coupling step is a call-site change inside the
Boris pusher's time-step loop; the collision physics does not have
to be written or validated again. The rough cost of Phase 2b is
reduced from ``write a new MCC module and validate it'' to
``integrate an existing validated MCC module''. That is a meaningful
cost reduction.

\paragraph{4. It isolates electron-kinetics effects from spatial
transport effects, which is the methodological prerequisite for
the non-local transport study.} The workplan §5.5 question cannot
be answered without first validating the local-field-approximation
kinetics layer on its own. The current deliverable performs that
validation. When Phase 2b runs the spatial PIC-MCC, any
disagreement between the spatial result and the cell-by-cell
surrogate can be cleanly attributed to genuine non-local transport,
because Tiers 1 and 2 and the 0D Tier 3 cross-check have already
been shown to agree with each other at the reference condition.
Without the current deliverable, Phase 2b would have to re-derive
all of those checks before it could interpret its own results.

% =====================================================================
\section{Remaining gap}
\label{sec:gap}

The current deliverable is explicit about what is \emph{not}
included. To avoid any ambiguity:

\begin{itemize}
  \item \textbf{No spatial PIC-MCC run.} The Tier 3 MCC collision
        core has not been coupled to the existing Boris pusher at
        M07/M08. No spatially resolved $k_j(r, z)$ maps have been
        extracted. No spatial steady-state has been computed.
  \item \textbf{No direct non-local transport validation in the
        TEL geometry.} The workplan §5.5 question about whether
        non-local transport moves the fluorine profile is
        unanswered in the current deliverable. The 0D MCC cross-
        check cannot answer it by construction.
  \item \textbf{No DTPM Picard-loop integration of
        \texttt{get\_rates\_pinn()}.} The Tier 2 API is built and
        validated against BOLSIG$^+$, but the Stage 10 DTPM Picard
        loop has not yet been modified to call it. The
        fluorine-profile comparison against Mettler 2025 has not
        yet been re-run with the Tier 2 surrogate in place of the
        Arrhenius block.
  \item \textbf{No Poisson self-consistency in the 0D MCC cross-
        check.} The current Tier 3 runs use a prescribed uniform
        field rather than a self-consistent space-charge field,
        which is a deliberate design choice for the 0D validation
        but will need revisiting in Phase 2b if the quasi-neutral
        approximation is found to be insufficient.
  \item \textbf{No secondary electrons in the Tier 3 MCC
        ionisation channel.} The ionisation scattering kernel
        loses threshold energy but does not create a secondary
        electron. At the sub-percent ionisation fractions of the
        workplan cases this is acceptable for the 0D cross-check;
        it will need to be added for self-consistent spatial PIC-
        MCC.
  \item \textbf{No extended-pressure Tier 1 sweep.} The Tier 1
        lookup is at a single pressure (10\,mTorr). Pressure
        dependence of the EEDF is expected to be weak in the DTPM
        operating regime per workplan §4.6 but has not been
        verified over the full 5--30\,mTorr range.
\end{itemize}

These gaps are not presented here as excuses. They are presented so
that the supervisor can read the reconciliation, the technical
report, the slides, and the README as a single consistent story:
\emph{the current deliverable is Phase 2a (the kinetics-isolation
step); Phase 2b (the spatial integration step) is the explicit next
package.}

% =====================================================================
\section{Next steps}
\label{sec:next}

To close the Phase 2 workplan in full, the recommended next steps
are:

\begin{enumerate}
  \item \textbf{Phase 2b: integrate the Tier 3 MCC collision core
        with the Boris pusher at M07/M08.} The MCC module at
        \texttt{tier3\_picmcc/mcc\_module.py} exposes a pure-Python
        \texttt{run\_mcc(...)} entry point. The integration is a
        call-site change inside the pusher's time-step loop. The
        main engineering work is: (i) replacing the prescribed
        uniform field with the pusher's cell-local $E_\theta(r, z,
        t)$; (ii) routing the EEDF accumulation into per-cell
        buffers on the masked TEL mesh; (iii) adding a
        post-processing wrapper on top of module M08 to extract
        spatially resolved $k_j(r, z)$. Estimated effort: 2--3
        weeks of software engineering plus 2--3 days of compute.
  \item \textbf{Run the three workplan cases with spatial
        resolution.} Same (power, pressure, $x_{\mathrm{Ar}}$)
        triples as the 0D cross-check in this deliverable, now on
        the TEL masked mesh. Extract $f_0(\varepsilon; r, z)$ and
        the volume-averaged $f_0^{\mathrm{avg}}(\varepsilon)$ for
        the ICP region. Compare against the Tier 1 lookup at the
        volume-averaged $E/N$; the agreement of these two
        quantities is the direct 0D comparison the workplan §5.4
        asks for.
  \item \textbf{Perform the spatial comparison.} Plot
        $k_{\mathrm{iz}}(r, z)$ from the spatial PIC-MCC alongside
        $k_{\mathrm{iz}}(r, z)$ from the cell-by-cell Tier 2
        surrogate lookup. If the spatial ionisation profile is
        broader than the local approximation predicts, that is the
        signature of non-local transport.
  \item \textbf{Couple the Tier 2 surrogate into the DTPM Picard
        loop.} This is independent of Phase 2b and can be done in
        parallel. Replace the current Arrhenius rate evaluation in
        the Stage 10 DTPM Picard loop with a batch call to
        \texttt{get\_rates\_pinn()}. Re-run the Mettler 2025
        fluorine-profile validation. If the centre-to-edge drop
        changes by less than 2 percentage points, the Maxwellian
        Arrhenius baseline stands as adequate for the fluorine
        profile. If it changes by more than that, Tier 2 becomes
        the production kinetics module.
  \item \textbf{Insert the PIC-MCC-derived $k_j(r, z)$ into the
        DTPM species transport solver.} Re-compute the
        [\text{F}] profile with the spatially-resolved MCC rates
        rather than the cell-by-cell surrogate lookup. Compare the
        centre-to-edge [\text{F}] drop against the Arrhenius
        baseline and the Tier 2 surrogate baseline. This is the
        workplan §5.4 Step 3.4.3 \emph{bottom-line test} for
        whether non-local transport matters.
  \item \textbf{(Optional) Extended-pressure Tier 1 sweep.} Repeat
        the BOLSIG$^+$ batch at 5 and 20\,mTorr. If the
        aggregated-rate pressure dependence is non-trivial,
        retrain the Tier 2 surrogate with pressure as a third
        input dimension.
  \item \textbf{(Optional) Add secondary-electron creation to the
        Tier 3 MCC ionisation channel.} Required for
        self-consistent spatial PIC-MCC if the ionisation fraction
        climbs above approximately 1\,\% during the run. The
        Opal-Peterson-Beaty energy-sharing model is the standard
        choice.
\end{enumerate}

Phase 2a (the current deliverable) and Phase 2b (steps 1--3 above)
together answer the full workplan. Steps 4--5 are orthogonal and
can proceed in parallel as long as someone other than the Tier 3
author is available for the DTPM Picard-loop integration work.

\vspace{1em}
\hrule
\vspace{0.5em}

\noindent
\textbf{Summary for the supervisor.} The current Phase 2 deliverable
is the electron-kinetics isolation step. Tiers 1 and 2 are complete
to their acceptance gates; Tier 3 is the collision core plus a 0D
cross-check, not the spatial PIC-MCC run the workplan ultimately
asks for. The gap is real, it is documented explicitly, and it is
structured so that the remaining work (Phase 2b) can begin as an
integration step on top of a validated collision core rather than
as a new module build. The request is that this deliverable be
accepted as Phase 2a and that Phase 2b be gated as a follow-on
package with its own time budget.

\end{document}
